% ============================================================
% FRONT MATTER -- OralGPT-Family Technical Report template
% Replace the title, authors, affiliations, abstract, date and links.
% Github  -> \ogsetlinkC      Hugging Face -> \ogsetlinkB
% ============================================================
\ogsettitle{OralGPT-Family Technical Report Template: \\ Title of Your Report Goes Here}

\ogsetauthors{Jing Hao$^{1}$, Jiahao Bao$^{1}$, Siyuan Dai$^{2}$,
              Yuci Liang$^{3}$, Xinpeng Lu$^{4}$}

\ogsetaffiliations{$^{1}$University of Hong Kong,
                   $^{2}$University of Pittsburgh,
                   $^{3}$Huazhong University of Science and Technology,
                   $^{4}$University of Science and Technology of China}

\ogsetteam{OralGPT-Family}

\ogsetemail{isjinghao@gmail.com}

\ogsetlinkA{https://github.com/isjinghao/OralGPT}
\ogsetlinkB{https://huggingface.co/OralGPT}
\ogsetlinkC{https://github.com/isjinghao/OralGPT}

\ogsetdate{2026-11-22}
\ogsetreporttype{Technical Report}
\ogsetstatus{OralGPT-Family Technical Report template}

\ogsetabstract{%
  This is the abstract placeholder of an OralGPT-Family technical report.
  Use the first one or two sentences to state the problem the report
  addresses and why it matters. Then summarise the proposed model, dataset,
  benchmark or system in a few sentences, naming its key components without
  going into implementation detail. Briefly describe how it is trained or
  constructed and how it is evaluated. Finish with the main findings, the
  released resources (models, data, code), and the intended use. A typical
  abstract is a single paragraph of 150--250 words. All text in this
  document is placeholder content that demonstrates the layout of the
  template; it does not describe a real system or real results.}

\ogmakefrontmatter
